\chapter{The open-source libraries, one by one}
\label{sec:libraries}

The method chapters told us what is worth running; this one asks who
actually ships it, and in what condition. Every version number,
maintenance signal, and feature claim in this chapter was checked against
the project's own repository, package index, release notes, or
documentation on August 22, 2026, and re-checked on August 23; the
summary appears in Table~\ref{tab:libraries}. Maintenance status is a fact with a date on it,
not a permanent verdict, and it matters here as much as algorithmic content:
several famous names below are no longer safe foundations. A useful way to
read the section is to sort the field by what each project fundamentally is:
a few are \emph{execution frameworks} that run and schedule trials, most are
\emph{search engines} that propose configurations and expect someone else to
run them, and two are \emph{services} wrapping a search engine behind a
network API. Our product needs exactly one good member of the first kind and
a curated set of the second.

\section{Execution frameworks}

\paragraph{Ray Tune.}
Part of Ray (release 2.57.0, August 11, 2026, Apache 2.0), Tune runs
distributed trials natively on a Ray cluster, with fault tolerance and a
stable Tuner interface. Its scheduler side is genuinely strong: ASHA,
Hyperband, median stopping, a BOHB scheduler, and the only mainstream
native distributed implementations of population-based training, PBT with
replay, and PB2, with the documented requirement that the trainable
support checkpoint save and restore \citep{liaw2018tune, moritz2018ray}.
Two findings from the release history temper the picture. First, across
the five Ray releases from December 2025 to August 2026, every
Tune-specific change was a documentation restructure, example, or bug fix;
no new algorithm or capability landed in that window. Second, the roster
of bundled search integrations has been deliberately pruned: Dragonfly,
SigOpt, scikit-optimize, and FLAML's searchers are gone from the current
documentation, leaving wrappers for Ax, bayesian-optimization, BOHB (whose
upstream is unmaintained, see below), HEBO, Hyperopt, Nevergrad, Optuna,
and ZOOpt. Core Tune has no native multi-objective support: the tuner
takes one metric and one mode, and Pareto workloads pass through the
Optuna searcher, which accepts lists of metrics. The fair summary: a
first-class execution and scheduling substrate whose own search
intelligence is coasting, which is precisely the shape of dependency we
want underneath a search layer of our own.

\paragraph{Syne Tune.}
Originally from AWS, now in its own community organization (release
0.16.0, August 3, 2026, Apache 2.0, about 420 GitHub stars), Syne Tune is
the closest existing thing to the product this document proposes: a broad
method zoo (multi-fidelity, multi-objective, transfer, and
population-based training all claimed in its README) behind a backend
abstraction that can execute trials locally, on SageMaker, or inside a
simulator that replays tabular and surrogate benchmarks faster than real
time \citep{salinas2022synetune}. That simulator is unique in the field
and exactly the right idea for cheap algorithm regression tests. The
concerns are social rather than technical: a small community and a
post-corporate maintenance arrangement, with no Ray backend, so adopting
it as our substrate would mean owning a second execution stack beside the
one RLlib already requires. We take its backend-abstraction design and its
simulator idea as blueprints rather than its code as foundation.

\paragraph{Katib.}
Kubeflow's Kubernetes-native tuner (v0.19.0, October 2025, Apache 2.0)
runs trials as Kubernetes workloads defined by custom resources, with
early stopping and a set of suggestion services that mostly wrap Optuna,
Goptuna, and Hyperopt. It is the right shape for an organization whose
platform is Kubeflow; on a Ray shop it duplicates the execution layer we
already have while adding no search capability we could not get from the
wrapped engines directly.

\paragraph{Determined and NNI.}
Two cautionary entries. Determined, the training platform whose adaptive
ASHA implementation was well regarded, has had no repository activity
since March 2025 and no release since then either, without any formal
notice; do not build on it. NNI, Microsoft's once-popular AutoML toolkit,
is formally archived (last release September 2023), and its fourteen
thousand GitHub stars are a trap for anyone ranking libraries by
popularity.

\section{Search engines}

\paragraph{Optuna.}
Optuna (4.9.0 in June 2026, 5.0 release candidate in August 2026, MIT) is
where sampler development is most active in 2026
\citep{akiba2019optuna}. The 5.0 line promotes a multivariate TPE with a
constant-liar strategy for parallel workers to the default sampler;
recent 4.x releases grew the GP sampler into constrained and
multi-objective territory (a log-EHVI acquisition, plus batch-proposal
strategies), an automatic sampler-selection module ships through the
OptunaHub plugin registry, and a Rust core accelerates the standard
samplers. Its define-by-run interface remains the best available way to
express conditional search spaces, and NSGA-II is the default for
multi-objective studies, with the multi-objective TPE variant
\citep{ozaki2020motpe} now folded into the standard sampler rather than
shipped separately. The limitation is the mirror image of Ray
Tune's: Optuna executes nothing. Its documented multi-node story is a
relational database that workers poll, with an explicit warning about
file-based storage over NFS, and it has no population-based training.
Optuna behind Ray Tune's OptunaSearch is the pairing both projects
document, and it is the pairing new tools reach for: DataRobot's syftr,
an agent-workflow optimizer released in 2025, states in its README that
it builds on Ray for distributed execution and Optuna for
define-by-run search and multi-objective optimization.

\paragraph{Ax and BoTorch.}
Meta's pair covers the model-based frontier: BoTorch (0.18.1, June 2026,
MIT) supplies the Gaussian-process machinery and the modern acquisition
functions, including the qLogNEHVI family that Chapter~\ref{sec:moo}
recommends, and its own documentation warns users off the older
non-log variants; Ax (1.3.1, June 2026, MIT) wraps it in a clean ask-tell
client aimed at practitioners, defaulting to noisy expected hypervolume
improvement for multi-objective studies
\citep{balandat2020botorch}. BoTorch's README is candid that end users
not doing BO research should use Ax instead. Neither executes trials;
both are exactly the acquisition engine our GP searchers should be built
on rather than reimplemented.

\paragraph{SMAC3.}
The AutoML group's random-forest Bayesian optimizer (2.4.0, April 2026,
BSD) remains actively developed, with multi-fidelity and multi-objective
support across its facades and a 2.4.0 change of default surrogate to
scikit-learn's random forest \citep{lindauer2022smac3}. It is the
canonical engine for large conditional spaces and the tuner used as the
sequential baseline in the BG-PBT paper. Its ergonomics are single-machine
and academic; we would wrap it for awkward spaces rather than adopt its
frame.

\paragraph{NePS.}
Neural Pipeline Search from the same group (0.16.0, March 2026, Apache
2.0, under a hundred stars but actively developed) is the reference home
of the prior-and-multi-fidelity methods we care about: PriorBand, ifBO,
and a multi-objective variant with expert priors, plus grammar-based
hierarchical architecture spaces. Its parallelism model is workers
coordinating through a shared results directory, simple and DDP-friendly,
though the same design that Optuna's documentation warns can race over
NFS. For us it is a source of reference implementations and a
correctness oracle for our own PriorBand and ifBO ports, not a substrate.

\paragraph{DEHB and HpBandSter.}
The two AutoML-group bandit engines are both frozen, in different ways.
DEHB (0.1.2, mid 2024) announces in its README that it is maintained for
stability but no longer actively developed; the algorithm remains the
evidence-backed recommendation for small-budget RL tuning
(Chapter~\ref{sec:workloads}), so the pragmatic path is to implement DEHB
against our own interfaces with the frozen package as a reference.
HpBandSter, the reference BOHB implementation, opens its README with
``not maintained anymore''; notably, Ray Tune's BOHB searcher still
depends on it.

\paragraph{Hyperopt, Nevergrad, FLAML.}
Hyperopt (0.3.0, July 2026) received its first release since 2021, a
compatibility revival with no new methods; it matters today mainly as the
legacy TPE many codebases still import. Nevergrad (1.0.12, April 2025,
sixteen months without a release at our check date) remains a broad
evolutionary toolbox whose NGOpt wizard picks among CMA-ES, differential
evolution, and dozens of other solvers by problem features; a good
non-BO baseline for bake-offs. FLAML (2.6.0, April 2026, active) is the
cost-frugal specialist: its CFO and BlendSearch searchers start cheap and
spend more only when justified, with explicit fields for the cost-related
hyperparameters \citep{wang2021flaml}; its Ray integration, however, has
disappeared from Ray's documentation, so using it means wiring it
ourselves.

\paragraph{OpenBox.}
An academic black-box optimization service (0.9.0, March 2026, MIT) with
an unusually broad claim sheet: multi-objective, constraints, transfer,
early stopping, and a REST service mode \citep{li2021openbox}. Activity
is modest and burst-like, and the service architecture duplicates
infrastructure a Ray shop already runs. Worth watching, not depending on.

\section{Services and partial offerings}

\paragraph{Open-source Vizier.}
Google's Vizier (client-server, Apache 2.0) publishes the most complete
description of an industrial default anywhere: a Mat\'ern-5/2 GP with
MAP-fit hyperparameters, an output-warping stack that tames outliers, a
deliberately explorative UCB coefficient, trust regions, and an
evolutionary acquisition maximizer, benchmarked in its own report against
algorithms accessed through Ray Tune \citep{song2022ossvizier,
song2024vizier}. The repository is active but the package on PyPI had not
been released for eighteen months at our check date. Its value to us is
the design disclosure more than the dependency.

\paragraph{Weights and Biases sweeps.}
The open pieces are the client and a small suggestion library (random,
grid, a Bayesian sampler, Hyperband stopping) whose standalone package
last shipped in 2022; orchestration and state live in the proprietary
service. Fine as a convenience for W\&B customers; not an engine to build
on.

\paragraph{New arrivals worth a footnote.}
The visible new energy of 2025 and 2026 sits in agentic tooling rather
than classical search: an official Optuna server exposing studies to
language-model agents, and a fast-growing project that has coding agents
iterate on training configurations directly. None of these publish
controlled evidence against the methods of
Chapters~\ref{sec:bayesian}--\ref{sec:transfer} at matched budgets, so
they change nothing in our plan, but they are worth rechecking in a year.

\begin{table}[t]
\centering
\small
\caption{Open-source tuning libraries as checked on August 22, 2026.
``Maintained'' reflects release and commit activity at that date, not a
prediction. Stars are approximate GitHub counts.}
\label{tab:libraries}
\begin{tabular}{@{}llll@{}}
\toprule
Library & Release (date) & Status & One-line reading \\
\midrule
Ray Tune & Ray 2.57.0 (2026-08) & active; Tune coasting & best executor; search layer aging \\
Optuna & 4.9.0; 5.0rc1 (2026-08) & very active & best samplers; no executor \\
Ax & 1.3.1 (2026-06) & active & polished BO client over BoTorch \\
BoTorch & 0.18.1 (2026-06) & active & the acquisition engine to reuse \\
Syne Tune & 0.16.0 (2026-08) & active, small team & broadest methods; design blueprint \\
SMAC3 & 2.4.0 (2026-04) & active & conditional-space BO plus racing \\
NePS & 0.16.0 (2026-03) & active, small team & PriorBand and ifBO reference code \\
DEHB & 0.1.2 (2024-07) & stability-only & algorithm good; package frozen \\
HpBandSter & 0.7.4 (2018-11) & unmaintained & BOHB reference; do not depend \\
Hyperopt & 0.3.0 (2026-07) & compatibility only & legacy TPE \\
Nevergrad & 1.0.12 (2025-04) & slowing & evolutionary baseline (NGOpt) \\
FLAML & 2.6.0 (2026-04) & active & cost-frugal searchers \\
OpenBox & 0.9.0 (2026-03) & sporadic & broad claims, small team \\
OSS Vizier & 0.1.24 (2025-02) & repo active, releases stalled & reference design disclosure \\
Katib & 0.19.0 (2025-10) & active & Kubernetes-native; redundant on Ray \\
Determined & 0.38.1 (2025-03) & dormant & no activity since early 2025 \\
NNI & 3.0 (2023-09) & archived & formally retired \\
Orion & 0.2.7 (2023-03) & dormant & no release in over three years \\
W\&B sweeps & lib 0.2.0 (2022-08) & minimal open core & orchestration is proprietary \\
\bottomrule
\end{tabular}
\end{table}
